\documentclass[12pt]{article}
\usepackage[T1]{fontenc}
\usepackage[utf8]{inputenc}
\usepackage{lmodern}
\usepackage{textcomp}
\usepackage{enumitem}
\usepackage{geometry}
\geometry{margin=1in}
\begin{document}
\begin{center}
Answer Key - Philosophy - Undergraduate Level Assessment 23 \
\small (Answers are ordered exactly as in the question paper.)
\end{center}

\section*{Multiple Choice}
\begin{enumerate}[label=\arabic*.]
\item \textbf{Answer:} A
\\ \textit{Options:} A) to minimize damage incurred by all individual living beings.; B) to promote the functional integrity of ecosystems.; C) to minimize the suffering of all sentient creatures.; D) to ensure the survival of endangered species.; E) to minimize damage to the ozone layer.
\\ \textit{Note:} Baxter's argument emphasizes the ethical responsibility to consider the well-being of all living beings in environmental decision-making, advocating for trade-offs that prioritize minimizing harm to ensure a balanced coexistence within ecosystems.
\item \textbf{Answer:} A
\\ \textit{Options:} A) the significance of moral intentions.; B) epistemological responsibility.; C) the role of virtue in morality.; D) the value of pleasure.; E) the importance of individual rights.
\\ \textit{Note:} The correct answer highlights that Ross critiques utilitarianism for its focus solely on the outcomes of actions rather than considering the moral intentions behind those actions, which he argues are crucial for true ethical evaluation.
\item \textbf{Answer:} B
\\ \textit{Options:} A) slippery slope; B) complex cause; C) red herring; D) ad hominem; E) hasty generalization
\\ \textit{Note:} The correct answer is "complex cause" because this fallacy occurs when an individual oversimplifies a situation by attributing an event's outcome to a single factor, ignoring the multiple interrelated causes that typically contribute to complex events.
\item \textbf{Answer:} D
\\ \textit{Options:} A) follow religious edicts.; B) follow historical precedents.; C) follow the teachings of philosophers.; D) follow Nature as our guide.; E) follow our conscience.
\\ \textit{Note:} Cicero posits that by following Nature as our guide, we align ourselves with the inherent moral order and rational principles that govern the universe, ensuring that our decisions are in harmony with true ethical values.
\item \textbf{Answer:} A
\\ \textit{Options:} A) "Man is a rational being."; B) "All bachelors are unmarried."; C) "Night follows day."; D) "The sum of the angles of a triangle is 180 degrees."; E) "The same thing cannot be affirmed and denied at the same time."
\\ \textit{Note:} The statement "Man is a rational being" is considered a self-evident proposition by Aquinas, as it reflects the intrinsic nature of humanity, while other options likely do not possess the same level of inherent clarity or universal acceptance.
\end{enumerate}

\section*{True / False}
\begin{enumerate}[label=\arabic*.]
\setcounter{enumi}{5}
\item \textbf{Answer:} True
\\ \textit{Options:} A) True; B) False
\\ \textit{Note:} The statement is true because it implies that the decision to raise tuition is contingent upon a recommendation from the board of trustees, which aligns with the governance structure of the university.
\item \textbf{Answer:} False
\\ \textit{Options:} A) True; B) False
\\ \textit{Note:} Research shows that the majority of individuals who experiment with illegal drugs do not transition to regular use in adulthood, as many people outgrow such behaviors or choose to abstain.
\item \textbf{Answer:} True
\\ \textit{Options:} A) True; B) False
\\ \textit{Note:} Paley asserts that the complexity and purposefulness observed in nature, which he describes as contrivance, signifies a divine creator because it is not a product of human craftsmanship but rather an intricate design inherent in the natural world.
\item \textbf{Answer:} True
\\ \textit{Options:} A) True; B) False
\\ \textit{Note:} The natural law fallacy relies on the erroneous assumption that because something occurs in nature, it is inherently moral, thereby misrepresenting moral reasoning through an inappropriate analogy between natural phenomena and ethical principles.
\item \textbf{Answer:} False
\\ \textit{Options:} A) True; B) False
\\ \textit{Note:} Plato contends that true beauty exists in the realm of ideal forms, which are abstract and not found in the physical world, thus making true beauty inaccessible in everyday objects.
\end{enumerate}

\section*{Long Form Response}
\begin{enumerate}[label=\arabic*.]
\setcounter{enumi}{10}
\item \textbf{Answer:} The fallacy of ignorance of refutation occurs when an individual claims that a proposition cannot be refuted due to a perceived or feigned inability to engage in the process of refutation. This fallacy leads to confusion, as it creates an illusion that the lack of counterarguments equates to the validity of the argument being presented. In reality, the inability to refute does not imply that the argument is sound; rather, it suggests a failure in critical engagement or understanding of the topic. Consequently, this fallacy undermines meaningful discourse by stifling debate and preventing the exploration of alternative perspectives, ultimately hindering the advancement of philosophical inquiry.
\\ \textit{Note:} This answer is correct because it accurately identifies the fallacy of ignorance of refutation as a misconception where the absence of counterarguments is mistakenly interpreted as evidence for the validity of an argument, thereby highlighting the importance of critical engagement in philosophical discourse.
\item \textbf{Answer:} The two principal perspectives on the meaning of life are the pessimist's view and the optimist's view, each offering distinct implications for human existence and well-being. The pessimist contends that life is inherently filled with suffering, futility, and disillusionment, suggesting that any search for meaning is ultimately futile. This perspective can lead to a sense of despair but also encourages a stark honesty about the human condition, prompting individuals to confront their limitations and the transient nature of existence. Conversely, the optimist sees life as imbued with potential, purpose, and the possibility of happiness, positing that individuals can create meaning through relationships, achievements, and personal growth. This viewpoint fosters resilience and hope, encouraging individuals to pursue fulfillment and well-being despite life's challenges. Together, these perspectives highlight the complex nature of human existence, underscoring the balance between acknowledging suffering and embracing hope.
\\ \textit{Note:} The answer correctly identifies the pessimist's view as one that emphasizes suffering and futility in life, highlighting its implications for honesty about human limitations, which contrasts with the optimist's perspective that seeks meaning and well-being.
\end{enumerate}

\vfill
\noindent\textit{End of Answer Key}
\end{document}